%%\setmainfont{FiraSans-Regular}[RawFeature={+lnum}]
%%\setmathfont[Extension = .otf]{FiraMath-Regular}
%%\setmathfont[version = bold, FakeBold=4, Extension = .otf]{FiraMath-Regular}
%
%
%\usepackage[sfdefault,lining]{FiraSans}
%\usepackage{firamath-otf}
%\setmathfont[version=bold,FakeBold=3, Extension=.otf]{FiraMath-Regular}
%\setboldmathrm{FiraSans}[
%	UprightFont=*-Bold,
%	BoldFont=*-ExtraBold,
%	ItalicFont=*-BoldItalic,
%	BoldItalicFont=*-ExtraBoldItalic,
%	Extension=.ttf]
%\usepackage{sfmath}               % wandelt Mathe-„Roman“ → Sans um
%\setmainfont{Fira Sans}      % Fließtext
%\setsansfont{Fira Sans}      % \textsf u.ä.
%\setmonofont{Fira Mono}      % Typewriter
%\setmathfont{Fira Math}      % alle Mathe-Glyphen
%
%%\usepackage[default,scale=1]{opensans}

\usepackage{iftex}
\ifXeTeX
  \usepackage{fontspec}
  \usepackage[math-style=ISO]{unicode-math}
  \setmainfont{Fira Sans}
  \setsansfont{Fira Sans}
  \setmonofont{Fira Mono}
  \setmathfont{Fira Math}
  \setmathfont[version=bold, FakeBold=3]{Fira Math}
\else\ifLuaTeX
  \usepackage{fontspec}
  \usepackage[math-style=ISO]{unicode-math}
  \setmainfont{Fira Sans}
  \setsansfont{Fira Sans}
  \setmonofont{Fira Mono}
  \setmathfont{Fira Math}
  \setmathfont[version=bold, FakeBold=3]{Fira Math}
\else
  % Fallback fuer pdfLaTeX
  \usepackage[sfdefault]{FiraSans}
  \usepackage{FiraMono}
\fi\fi